% The companion-paper note. Every paper in the series carries it so a reader who arrives at one
% knows what the other two contain and, more importantly, which claims are NOT in the paper they
% are holding. Three papers were carved out of one 46-page draft; the carve is documented in
% docs/papers/README.md, and the split's completeness is machine-checked by
% scripts/verify_split_coverage.py against the pre-split source.

\paragraph{This paper is one of three.}
One study produced three separable results, and they are reported separately because they answer
different questions, rest on different experiments, and are useful to different readers. All three
share one harness, one artifact release and one corrections register.
\begin{description}\itemsep0pt
  \item[I. The conditional lever.] Whether optimizing a self-refine loop's verifier beats
    optimizing its generation prompt, and the quantity that decides it: $p_0$, the base rate of
    first drafts already satisfying every requirement.
  \item[II. What a checker cannot buy is the repair.] The loop's own ceiling, with no optimizer
    anywhere in it: how good the intrinsic critique is as a detector, what happens when the
    detector is made perfect, and which fix-side levers actually pay.
  \item[III. Your loop optimizer is not optimizing your loop.] A silent failure mode in
    componentized prompt optimization: pointing an optimizer at a component is not the same as
    optimizing it, the difference does not appear in held-out scores, and one integer per component
    catches it.
\end{description}
Shared artifact, including every cassette, pre-registration and the corrections register:
\texttt{github.com/loopgepa/whystuck}.
